\subsection{GP-1965-08: a rational coordinate field in the first Weyl division ring}
\label{dev:weyl-coordinate}

\paragraph{Question, provenance and scope.}
The definite question in occurrence \hyperref[item:1967-145.3]{1967-145.3},
repeating 1965-80.10, is to describe the automorphisms of the fraction division
ring of the first Weyl algebra. Work throughout over \(\mathbb C\), with
\[
 A_1=\mathbb C\langle p,q\rangle/(pq-qp-1),\qquad
 D_1=\operatorname{Frac}A_1=\mathbb C(q)(p;\partial_q).
\]
All automorphisms fix \(\mathbb C\) pointwise. The notebook contrasts the
known polynomial automorphisms with the unknown division-ring case. It does
not assert that every division-ring automorphism preserves
\(K=\mathbb C(q)\), and it does not formulate the endomorphism version of the
Dixmier conjecture. Those are different questions.

The bounded problem addressed here is to describe
\[
 N_K=\{\phi\in\operatorname{Aut}_{\mathbb C}D_1:\phi(K)=K\}
 \quad\hbox{and}\quad
 U_q=\{\phi\in\operatorname{Aut}_{\mathbb C}D_1:\phi(q)=q\}.
\]
The equality \(\phi(K)=K\), rather than mere containment, matters. The
strongest result proved below is an exact description of both groups and of
their intersection with \(\operatorname{Aut}A_1\). This is a complete
verification of a known special case, not a new solution of the full
classification problem. The normalizer formula is explicitly recorded by
\citet[\S5.3.3, final Remark, printed p.~66]{dumas2011invariants}, who refers
to \citet{alevdumas1998invariants}. The formal-series embedding used below is
in \citet[\S5.2.2, pp.~59--60]{dumas2011invariants}. The proofs are included
to make the reduction independently checkable; the marker \ours\ identifies
work re-derived here, without a novelty claim.

\paragraph{What the literature does and does not supply.}
\citet[Theorem~8.10]{dixmier1968weyl} describes the polynomial automorphism
group of \(A_1\). Passing to fractions extends every such automorphism, but
also admits rational shears and inner automorphisms that need not preserve
the polynomial algebra. Dumas separates actions extending from the algebra
from general actions on its fraction field; his printed p.~66 gives the
exact coordinate-field stabilizer while explicitly leaving the full group
undetermined. The current-status audit accompanying occurrence 1967-145.3
does not identify a later complete classification. The obstruction below
explains the logical gap in this particular argument; it is not offered as
a historical explanation of why other mathematicians have not solved the
problem.

\begin{lemma}[The coordinate centralizer]\ours\ \label{lem:coordinate-centralizer}
The centralizer of \(q\) in \(D_1\) is exactly \(K=\mathbb C(q)\).
\end{lemma}
\begin{proof}
Here is a formal construction that also justifies the use of negative
powers of \(p\). Let \(\mathcal S\) consist of series
\(T=\sum_{j\leq m}a_j(q)p^j\), with \(m\in\mathbb Z\) and
\(a_j\in K\). Exponents are bounded above and can be unbounded below.
Define multiplication by the generalized Leibniz rule
\[
 (ap^i)(bp^j)=\sum_{k\geq0}\binom{i}{k}a\,\partial_q^k(b)\,p^{i+j-k},
 \qquad i,j\in\mathbb Z.
\]
For a fixed output exponent only finitely many terms contribute: if the
input orders are at most \(m,n\), a contribution to exponent \(\ell\)
requires \(i+j\geq\ell\), hence \(i\geq\ell-n\) and
\(j\geq\ell-m\). Thus there is no analytic convergence assertion.

For completeness, associativity follows coefficient by coefficient from
Leibniz' rule and the identity
\[
 \sum_{t=0}^{v}\binom jt\binom i{u+v-t}\binom{u+v-t}u
 =\binom iu\binom{i+j-u}v
 \qquad(u,v\geq0).
\]
Indeed, both sides are the coefficient of
\(a\,\partial_q^u(b)\,\partial_q^v(c)\) in the two products of
\(ap^i,bp^j,cp^k\). The identity is the generalized Vandermonde identity
after factoring out \(\binom iu\), and holds for negative as well as
nonnegative integers \(i,j\). Finite coefficient sums justify its use for
series.

Every nonzero series is invertible. If its leading term is \(a_mp^m\),
factor it as \(a_mp^m(1+R)\), where \(R\) has order at most \(-1\).
The series \(\sum_{k\geq0}(-R)^k\) is well defined coefficientwise and is
a two-sided inverse of \(1+R\). The monomial inverse is
\(p^{-m}a_m^{-1}\). Consequently \(\mathcal S\) is a division ring.
The rule \(pf=fp+f'\) embeds \(K[p;\partial_q]\) into \(\mathcal S\);
uniqueness of the polynomial expansion makes this embedding injective.
Its Ore fraction division ring therefore embeds as well.

Since \(q'=1\) and \(q^{(k)}=0\) for \(k\geq2\), the same rule gives
\[
 [p^j,q]=jp^{j-1}\quad\hbox{for every }j\in\mathbb Z,
 \qquad
 [T,q]=\sum_{j\leq m}j a_j(q)p^{j-1}.
\]
The exponents \(j-1\) are distinct. In characteristic zero the last
series vanishes only if \(a_j=0\) for every \(j\ne0\). Thus an element
of \(D_1\) commuting with \(q\) lies in \(K\). The reverse inclusion
is immediate.
\end{proof}

\begin{proposition}[All automorphisms fixing the coordinate]\ours\ 
\label{prop:rational-shear}
The map \(r\mapsto u_r\), where
\[
 u_r(q)=q,\qquad u_r(p)=p+r(q),\qquad r\in K,
\]
identifies the additive group of \(K\) with \(U_q\).
\end{proposition}
\begin{proof}
If \(\phi(q)=q\), then
\([\phi(p)-p,q]=\phi([p,q])-[p,q]=0\). The preceding lemma gives
\(\phi(p)=p+r(q)\). Conversely, replacing \(p\) by \(p+r\) and fixing
\(K\) preserves the relation \(pf-fp=f'\). It defines an automorphism of
\(K[p;\partial_q]\), with inverse obtained by replacing \(r\) by
\(-r\), and hence an automorphism of its fraction division ring.
Finally \(u_ru_s=u_{r+s}\); distinct functions have distinct images of
\(p\).
\end{proof}

\begin{lemma}[The automorphisms of the coefficient field]\ours\ 
Every \(\mathbb C\)-automorphism of \(K\) sends \(q\) to
\[
 \theta(q)=\frac{aq+b}{cq+d},\qquad ad-bc\ne0.
\]
\end{lemma}
\begin{proof}
Let its image be a nonconstant rational function \(\theta\), with
rational inverse substitution \(\psi\). A nonconstant rational function
defines a morphism of \(\mathbb P^1\) of degree equal to the common degree
of a coprime pair of homogeneous defining polynomials. Degrees multiply
under composition: composing the two homogeneous pairs gives the product
degree and has no common projective zero, because neither original pair
has one. Since \(\psi\circ\theta\) is the identity, both degrees are one.
The degree-one expressions are precisely the displayed fractions with
nonzero determinant. Conversely such a fraction has a degree-one rational
inverse and defines a field automorphism.
\end{proof}

\begin{proposition}[Normalizer of a rational coordinate field]\ours\ 
\label{prop:coordinate-normalizer}
The elements of \(N_K\) are exactly
\[
 \phi_{\theta,r}(q)=\theta(q),\qquad
 \phi_{\theta,r}(p)=\frac{1}{\theta'(q)}p+r(q),
 \quad \theta(q)=\frac{aq+b}{cq+d},\quad ad-bc\ne0,\quad r\in K.
\]
Restriction gives a split exact sequence
\[
 1\longrightarrow (K,+)\longrightarrow N_K
 \longrightarrow\operatorname{Aut}_{\mathbb C}K\longrightarrow1.
\]
With \(\sigma_\theta=\phi_{\theta,0}\), its action on rational shears is
\(\sigma_\theta u_r\sigma_\theta^{-1}=u_{\theta'\,(r\circ\theta)}\).
\end{proposition}
\begin{proof}
Restriction to \(K\) is a field automorphism, so the preceding lemma
determines \(\phi(q)=\theta(q)\). Because \(\theta'\ne0\),
\([\theta'^{-1}p,\theta]=1\). Hence
\(\phi(p)-\theta'^{-1}p\) commutes with \(\theta(q)\). Since
\(\mathbb C(\theta(q))=K\), commuting with \(\theta(q)\) is equivalent
to commuting with \(q\), and Lemma~\ref{lem:coordinate-centralizer}
gives the required rational remainder \(r\).

Conversely, substitution on \(K\), together with
\(p\mapsto\theta'^{-1}p+r\), respects the Ore relation: for \(f\in K\),
\[
 [\theta'^{-1}p+r,f(\theta(q))]
   =\theta'^{-1}\partial_q(f(\theta(q)))=f'(\theta(q)).
\]
The image of a nonzero polynomial of order \(m\) has nonzero leading
coefficient \(f_m(\theta)\theta'^{-m}\), so the resulting homomorphism
of the Ore polynomial ring is injective and extends to its fractions.
It is onto: \(q=\theta^{-1}(\phi(q))\) lies in the generated division
ring, and then \(p=\theta'(q)(\phi(p)-r(q))\) does too.

Writing \(\psi=\theta^{-1}\), an explicit inverse is
\[
 \phi_{\theta,r}^{-1}(q)=\psi(q),\qquad
 \phi_{\theta,r}^{-1}(p)=\frac{1}{\psi'(q)}p
                       -\frac{r(\psi(q))}{\psi'(q)}.
\]
The order of the coefficient and \(p\) in these expressions is essential.
The chain rule checks both compositions. The same chain rule shows that
the coefficient-field automorphism lifts with \(r=0\) give a section.
Conjugating \(u_r\) by this section yields
\(p\mapsto p+\theta'(q)r(\theta(q))\), proving the last formula.
The exact sequence uses composition of field automorphisms; if rational
maps are written as functions, substitution reverses their displayed
composition order. This convention prevents an accidental sign or order
error in a semidirect-product formula.
\end{proof}

\begin{corollary}[The polynomial boundary]\ours\ \label{cor:polynomial-boundary}
An automorphism \(\phi_{\theta,r}\) in
Proposition~\ref{prop:coordinate-normalizer} restricts to an automorphism
of \(A_1\) if and only if
\[
 \theta(q)=aq+b\quad(a\ne0),\qquad r(q)\in\mathbb C[q].
\]
In particular, a shear fixing \(q\) preserves \(A_1\) exactly when its
shear function is polynomial.
\end{corollary}
\begin{proof}
The PBW expansion in \(A_1=\mathbb C[q][p;\partial_q]\), together with
the unique formal expansion above, gives
\(A_1\cap K=\mathbb C[q]\). If \(\phi\) and \(\phi^{-1}\) preserve
\(A_1\), both \(\theta(q)\) and \(\theta^{-1}(q)\) are polynomials.
Their degrees multiply to one, so \(\theta\) is affine. Its derivative
is the nonzero constant \(a\). Subtracting \(a^{-1}p\) from
\(\phi(p)\) then shows that \(r\in A_1\cap K=\mathbb C[q]\).
Conversely affine substitution and a polynomial shear have polynomial
inverse formulas as displayed above, so they preserve \(A_1\) in both
directions.
\end{proof}

\begin{example}[A denominator changes the ambient problem]\ours\ 
The rational shear \(q\mapsto q,\ p\mapsto p+1/q\) is an automorphism
of \(D_1\), with inverse \(p\mapsto p-1/q\), and does not preserve
\(A_1\). Its hypotheses are explicit: \(q\ne0\) as an element of the
division ring, so \(q^{-1}\) exists there; the PBW intersection above
shows that it is not an element of the polynomial algebra. Similarly
\(q\mapsto-1/q,\ p\mapsto q^2p\) is a normalizer automorphism.
These examples exhibit the precise limitation of transferring the
polynomial generation theorem to the division ring.
\end{example}

\begin{remark}[Independent computation]\oursComputed\ 
The accompanying \path{code/verify_weyl.py} performs exact symbolic
checks, with no floating-point tolerances. It tests the negative-power
commutator identity in the formal Leibniz algebra, verifies both orders of
the inverse formula for several affine and nonaffine M\"obius maps and
several polynomial and rational shears, and checks the commutator by
applying the resulting first-order differential operators to independent
rational test functions. The output gives the parameter families and
number of checks. These finite checks detect ordering and sign errors;
they are not a substitute for the coefficientwise and chain-rule proofs
above. A second independent algebraic review is recorded in
\path{audit/proof_review.md}; its findings are incorporated into this
version.
\end{remark}

\paragraph{The obstruction still present.}
For a general \(\phi\in\operatorname{Aut}D_1\), put
\(Q=\phi(q)\), \(P=\phi(p)\). Necessarily \([P,Q]=1\) and the
division subring generated by \(P,Q\) is \(D_1\). The image
\(\mathbb C(Q)\) is self-centralizing, because centralizers are transported
by automorphisms. Our proof classifies the possibilities only after
\(\mathbb C(Q)=\mathbb C(q)\) is known. In general there is no reason
for that equality. The full problem requires controlling the orbit of
rational coordinate subfields, not merely their stabilizer. Moreover,
finding two elements with commutator one alone does not certify that they
generate the division ring; that further assertion cannot be inserted
without proof.

\paragraph{A precise next bounded task.}
One further research run could fix an explicit pair of rational
expressions \(P,Q\), verify \([P,Q]=1\), supply inverse expressions for
\(p,q\) to certify generation, and then determine whether its coordinate
field can be carried to \(\mathbb C(q)\) by a stated finite composition
of rational shears, Fourier exchange and inner conjugations. A finite
degree or denominator bound would make such a computation reproducible.
Success for that bounded family would give a factorization certificate
for those examples. It would not prove that the chosen transformations
generate every automorphism, nor convert a finite search with no
factorization into a counterexample. This is the exact limit of the
development undertaken here.
